\chaptertitle{PS2 · 위험태도, 확실성등가와 위험프리미엄}{공식 Exercise 1--5 · CARA/CRRA, 기대효용, wealth effect, 확실성등가}

\begin{sourcebox}
\textbf{이 장의 공식 출처}: MBF\_PS2\_2024-2.pdf, MBF\_PS2\_Solution\_2024-2.pdf. 공식 문항의 인물·상황·금액을 그대로 사용했다.
\end{sourcebox}

\exercise{PS2 Exercise 1}{공식 문제}
\begin{officialbox}
Payoff A는 6면체 주사위를 한 번 던져 나온 눈만큼의 달러를 받는 복권이고, Payoff B는 확실하게 3달러를 받는 선택이다.
\begin{enumerate}[label=(\alph*)]
\item 초기부가 0인 Binh의 효용함수는 모든 $W>0$에 대해
\[
-\frac{WU''(W)}{U'(W)}=1
\]
을 만족한다. A와 B 중 무엇을 선택하는가?
\item Guillaume의 효용함수는
\[
-\frac{u''(W)}{u'(W)}=k,\qquad k>0
\]
을 만족한다. 초기부 $Y_0=1$일 때 B를 선택한다는 사실을 알고 있다. 초기부가 충분히 커지면 선택을 바꾸는가? 바꾼다면 임계값을 구하고, 바꾸지 않는다면 그 이유를 설명하라.
\end{enumerate}
\end{officialbox}

\begin{strategybox}
주어진 식은 각각 상대위험회피도(RRA)와 절대위험회피도(ARA)다. 먼저 미분방정식에서 대표 효용함수를 복원한 뒤 기대효용을 비교한다.
\end{strategybox}

\begin{solutionbox}
\step{(a) RRA가 1이면 로그효용}
\[
-\frac{WU''(W)}{U'(W)}=1
\quad\Longrightarrow\quad
\frac{U''(W)}{U'(W)}=-\frac1W.
\]
양변을 적분하면
\[
\ln U'(W)=-\ln W+C
\quad\Longrightarrow\quad U'(W)=\frac{A}{W}.
\]
다시 적분하면 $U(W)=A\ln W+B$다. 양의 affine transformation은 선호를 바꾸지 않으므로 $U(W)=\ln W$로 둔다.

A의 기대효용은
\[
V_A=\frac16\sum_{j=1}^6\ln j
=\frac16\ln(1\cdot2\cdot3\cdot4\cdot5\cdot6)
=\ln\left(720^{1/6}\right).
\]
B의 효용은
\[
V_B=\ln3.
\]
비교하면
\[
V_A<V_B
\quad\Longleftrightarrow\quad
720^{1/6}<3
\quad\Longleftrightarrow\quad
720<3^6=729.
\]
따라서 B를 선택한다.

\step{(b) ARA가 일정하면 CARA 효용}
\[
-\frac{u''(W)}{u'(W)}=k
\quad\Longrightarrow\quad
u'(W)=Ce^{-kW}.
\]
대표 효용함수는
\[
u(W)=-e^{-kW}
\]
로 둘 수 있다. 초기부가 $Y_0$일 때 A의 기대효용은
\[
EU_A(Y_0)=-\frac16\sum_{j=1}^6 e^{-k(Y_0+j)}
=-e^{-kY_0}\left(\frac16\sum_{j=1}^6e^{-kj}\right).
\]
B의 효용은
\[
U_B(Y_0)=-e^{-k(Y_0+3)}=-e^{-kY_0}e^{-3k}.
\]
두 식의 공통 양의 인자 $e^{-kY_0}$는 비교에서 소거된다. 따라서 A와 B의 순위는 $Y_0$와 무관하다. $Y_0=1$에서 B를 선택했다면 모든 초기부에서 B를 선택한다.
\end{solutionbox}

\begin{answerbox}
\begin{enumerate}[label=(\alph*)]
\item $U(W)=\ln W$로 나타낼 수 있고, $720<729$이므로 \textbf{Payoff B}를 선택한다.
\item $u(W)=-e^{-kW}$인 CARA 선호에서는 초기부가 공통인자로 소거된다. 따라서 \textbf{선택을 바꾸지 않으며 임계 초기부도 없다}.
\end{enumerate}
\end{answerbox}

\begin{warningbox}
CARA의 ``wealth effect가 없다''는 말은 위험자산에 투자하는 \textbf{금액} 또는 주어진 복권의 선택이 초기부와 무관하다는 뜻이다. 모든 경제행동이 부와 무관하다는 뜻은 아니다.
\end{warningbox}

\exercise{PS2 Exercise 2}{공식 문제}
\begin{officialbox}
Raluca와 Yuki는 모두
\[
-\frac{WU''(W)}{U'(W)}=\frac12
\]
을 만족하는 선호와 초기부 1달러를 갖는다.
\begin{enumerate}[label=(\alph*)]
\item Argentina와 France가 각각 $1/2$ 확률로 이긴다고 생각한다. ``Argentina 승리 시 Raluca가 Yuki에게서 1달러를 받고, France 승리 시 반대로 지급''하는 내기에 두 사람 모두 동의하는가?
\item Raluca는 Romania 승리확률을 $3/4$, Yuki는 Japan 승리확률을 $3/4$로 믿는다. 동일한 구조의 내기에 두 사람 모두 동의하는가?
\end{enumerate}
\end{officialbox}

\begin{strategybox}
RRA $\gamma=1/2$인 CRRA 효용은 $U(W)=2\sqrt W$로 둘 수 있다. 내기를 수락하면 각 사람의 최종부는 2 또는 0, 거절하면 1이다.
\end{strategybox}

\begin{solutionbox}
\step{효용함수 복원}
CRRA 대표효용은
\[
U(W)=\frac{W^{1-\gamma}}{1-\gamma}.
\]
$\gamma=1/2$를 넣으면
\[
U(W)=2\sqrt W.
\]

\step{(a) 객관적 50--50 내기}
Raluca가 내기를 수락할 때 기대효용은
\[
\frac12U(2)+\frac12U(0)
=\frac12(2\sqrt2)+\frac12(0)=\sqrt2.
\]
거절하면
\[
U(1)=2.
\]
$\sqrt2<2$이므로 Raluca는 거절한다. Yuki에게도 완전히 대칭이므로 Yuki도 거절한다.

\step{(b) 서로 다른 주관적 믿음}
Raluca는 자신이 2달러가 될 확률을 $3/4$로 본다.
\[
EU_R=\frac34U(2)+\frac14U(0)=\frac32\sqrt2.
\]
이를 $U(1)=2$와 비교하면
\[
\frac32\sqrt2\ge2
\quad\Longleftrightarrow\quad
\sqrt2\ge\frac43,
\]
이는 참이다. 따라서 Raluca는 수락한다. Yuki 역시 자신의 관점에서 승리확률을 $3/4$로 믿으므로 같은 계산으로 수락한다.
\end{solutionbox}

\begin{answerbox}
\begin{enumerate}[label=(\alph*)]
\item 둘 다 거절한다. 공정한 평균보존위험을 위험회피자가 선호하지 않기 때문이다.
\item 둘 다 수락한다. 서로 다른 낙관적 믿음 때문에 각자 자신의 승리확률을 $3/4$로 평가한다.
\end{enumerate}
\end{answerbox}

\begin{warningbox}
(b)에서 실제로 양쪽이 동시에 $3/4$ 확률로 이길 수 있다는 뜻이 아니다. \textbf{주관적 확률이 서로 다르기 때문에} 거래가 성립한다는 문제다.
\end{warningbox}

\exercise{PS2 Exercise 3}{공식 문제}
\begin{officialbox}
Raluca의 선호는 모든 $W>0$에 대해
\[
-\frac{U''(W)}{U'(W)}=1
\]
을 만족하고 초기부는 1달러다.
\begin{enumerate}[label=(\alph*)]
\item 두 후보의 승리확률이 각각 $1/2$라고 생각할 때, 이기면 1달러를 받고 지면 1달러를 지급하는 내기를 수락하는가?
\item 다른 선거에서 West가 이기면 1달러를 받고 Trump가 이기면 1달러를 지급하는 내기에 무차별이다. Raluca가 믿는 West의 승리확률을 구하라.
\end{enumerate}
\end{officialbox}

\begin{strategybox}
ARA가 1인 CARA 효용은 $U(W)=-e^{-W}$다. (a)는 공정한 위험의 거절, (b)는 기대효용과 확실한 부 1의 효용을 같게 두는 문제다.
\end{strategybox}

\begin{solutionbox}
\step{(a)}
내기를 수락하면 최종부는 $1/2$ 확률로 2, $1/2$ 확률로 0이다. 기대부는 1로서 거절할 때의 확실한 부 1과 같다. 그러나 $U''<0$인 위험회피자에게
\[
\frac12U(2)+\frac12U(0)<U\left(\frac12\cdot2+\frac12\cdot0\right)=U(1).
\]
따라서 거절한다.

\step{(b)}
West 승리확률을 $\pi$라 하자. 무차별조건은
\[
\pi U(2)+(1-\pi)U(0)=U(1).
\]
$U(W)=-e^{-W}$를 넣으면
\[
-\pi e^{-2}-(1-\pi)=-e^{-1}.
\]
정리하면
\[
\pi(1-e^{-2})=1-e^{-1},
\]
\[
\pi=\frac{1-e^{-1}}{1-e^{-2}}
=\frac{e}{1+e}\approx0.731.
\]
\end{solutionbox}

\begin{answerbox}
\[
\text{(a) 거절},\qquad
\text{(b) }\boxed{\pi=\frac{e}{1+e}\approx73.1\%}.
\]
\end{answerbox}

\begin{warningbox}
(a)에서 ``기대값이 같으므로 무차별''이라고 쓰면 위험중립자의 답이다. 위험회피에서는 Jensen 부등식으로 확실한 1을 더 선호한다.
\end{warningbox}

\exercise{PS2 Exercise 4}{공식 문제}
\begin{officialbox}
게임쇼에서 참가자는 500,000달러를 받고 그만두거나, 네 선택지 중 하나에 답한다. 정답이면 1,000,000달러, 오답이면 32,000달러를 받는다.
\begin{enumerate}[label=(\alph*)]
\item 초기부 0인 Liqian은 RRA가 1이다. 그녀는 현재 500,000달러 제안을 거절하고 답하지만, 확실히 650,000달러를 받을 수 있었다면 그만두었을 것이라고 말한다. 정답확률에 대한 그녀의 믿음을 구하고, 왜 답을 선택하는지 설명하라.
\item 초기부 0인 Dario는 위험선호자이며 네 답이 모두 똑같이 가능하다고 생각한다. 그가 답할지 그만둘지 예측할 수 있는가?
\end{enumerate}
\end{officialbox}

\begin{strategybox}
Liqian은 로그효용을 가지며 650,000달러가 답하기 복권의 확실성등가다. Dario는 단순히 ``위험선호''라는 곡률 정보만 주어졌으므로 효용함수의 정확한 굽힘 정도가 필요하다.
\end{strategybox}

\begin{solutionbox}
\step{(a) Liqian의 주관적 정답확률}
$X_1=1{,}000{,}000$, $X_2=32{,}000$, 확실성등가를 $\widehat K=650{,}000$이라 하자. 정답확률을 $p$라 하면 로그효용의 무차별조건은
\[
p\ln X_1+(1-p)\ln X_2=\ln\widehat K.
\]
이를 풀면
\[
p=\frac{\ln(\widehat K/X_2)}{\ln(X_1/X_2)}
=\frac{\ln(650{,}000/32{,}000)}{\ln(1{,}000{,}000/32{,}000)}
\approx0.875.
\]
답하기 복권의 확실성등가는 650,000달러인데 실제 walk-away 금액은 500,000달러이므로
\[
CE(\text{답하기})=650{,}000>500{,}000.
\]
따라서 답한다.

\step{(b) Dario}
아무 정보가 없으므로 정답확률은 $1/4$이고 오답확률은 $3/4$다. 답하기의 기대금액은
\[
\frac14(1{,}000{,}000)+\frac34(32{,}000)=274{,}000,
\]
500,000달러보다 작다. 위험중립에 가까운 약한 위험선호자라면 그만둘 수 있다. 반면 효용함수가 충분히 강하게 볼록하면 1,000,000달러의 상방을 크게 평가하여 답할 수 있다. $u''(W)>0$만으로는 두 효용을 비교할 수 없다.
\end{solutionbox}

\begin{answerbox}
\[
\text{(a) }p\approx0.875\;(87.5\%),\qquad \text{Liqian은 답한다.}
\]
\[
\text{(b) 판단 불가. 위험선호의 정도, 즉 효용함수의 정확한 곡률이 더 필요하다.}
\]
\end{answerbox}

\begin{warningbox}
위험선호라고 해서 기대금액이 더 낮은 모든 복권을 반드시 고르는 것은 아니다. ``볼록하다''는 방향만 알려 줄 뿐, 두 점 사이의 효용차 크기까지 정하지 않는다.
\end{warningbox}

\exercise{PS2 Exercise 5}{공식 문제}
\begin{officialbox}
Deal or No Deal에서 상자에는 10달러 또는 2달러가 각각 $1/2$ 확률로 들어 있다.
\begin{enumerate}[label=(\alph*)]
\item 초기부 1이고 RRA가 1인 Valérie에게 은행가가 상자를 5달러에 사겠다고 한다. 제안을 수락하는가? 효용곡선으로 설명하라.
\item 초기부가 1보다 크면 결정을 바꾸는가? 바뀌는 임계 초기부를 구하라.
\item 초기부 0이고 RRA가 2인 Pavel에게 은행가가 $B$를 제안한다. Pavel이 수락하는 최소 $B$를 구하라.
\item Pavel의 상자에 대한 위험프리미엄 $\Pi$를 구하라.
\end{enumerate}
\end{officialbox}

\begin{strategybox}
Valérie는 로그효용, Pavel은 $U(W)=-1/W$로 나타낼 수 있다. ``Deal''과 ``No deal''의 최종부를 정확히 비교하고, 최소 제안액은 확실성등가이며 위험프리미엄은 기대값에서 확실성등가를 뺀 값이다.
\end{strategybox}

\begin{solutionbox}
\step{(a) Valérie}
초기부를 $Y=1$, 제안액을 $B=5$, 상자 금액을 $x_1=10$, $x_2=2$라 하자. Deal의 최종부는 $Y+B=6$이고 No deal의 최종부는 11 또는 3이다.
\[
\ln6\;\mathop{\gtrless}_{\text{No deal}}^{\text{Deal}}\;
\frac12\ln11+\frac12\ln3.
\]
로그의 단조성을 이용하면
\[
\ln6>\frac12\ln(33)
\quad\Longleftrightarrow\quad 36>33.
\]
따라서 Deal을 수락한다. 오목한 효용곡선에서 확실한 6의 효용이 11과 3의 기대효용보다 높다.

\step{(b) 초기부의 임계값}
일반적인 $Y$에서 Deal을 선택할 조건은
\[
\ln(Y+B)\ge\frac12\ln(Y+x_1)+\frac12\ln(Y+x_2).
\]
지수화하면
\[
(Y+B)^2\ge(Y+x_1)(Y+x_2).
\]
정리하면
\[
[2B-(x_1+x_2)]Y+B^2-x_1x_2\ge0.
\]
현재 수치를 넣으면
\[
-2Y+5\ge0
\quad\Longleftrightarrow\quad Y\le\frac52.
\]
따라서 $Y=5/2$에서 무차별이고, $Y>5/2$이면 No deal을 택한다. CRRA 효용에서는 초기부가 증가하면 같은 절대금액 복권의 상대적 크기가 작아져 행동이 달라질 수 있다.

\step{(c) Pavel의 최소 제안액}
RRA가 2인 CRRA 효용은
\[
U(W)=\frac{W^{1-2}}{1-2}=-\frac1W
\]
로 둘 수 있다. 초기부가 0이므로 수락조건은
\[
-\frac1B\ge\frac12\left(-\frac1{10}\right)+\frac12\left(-\frac12\right).
\]
오른쪽은 $-3/10$이므로
\[
\frac1B\le\frac3{10}
\quad\Longrightarrow\quad
\boxed{B\ge\frac{10}{3}}.
\]
최소 제안액, 즉 확실성등가는 $\widehat B=10/3$이다.

\step{(d) 위험프리미엄}
상자의 기대금액은
\[
\E[\widetilde x]=\frac12(10)+\frac12(2)=6.
\]
따라서
\[
\Pi=\E[\widetilde x]-CE
=6-\frac{10}{3}=\frac83.
\]
\end{solutionbox}

\begin{answerbox}
\begin{align*}
\text{(a)}\quad &\text{Valérie는 5달러 제안을 수락한다.}\\
\text{(b)}\quad &Y^*=\frac52;\;Y<\frac52\text{이면 Deal, }Y>\frac52\text{이면 No deal}.\\
\text{(c)}\quad &B_{\min}=\frac{10}{3}.\\
\text{(d)}\quad &\Pi=\frac83.
\end{align*}
\end{answerbox}

\begin{warningbox}
위험프리미엄은 $CE-\E[X]$가 아니라 \textbf{$\E[X]-CE$}다. 위험회피자라면 일반적으로 $CE<\E[X]$이므로 위험프리미엄은 양수다.
\end{warningbox}
